\chapter{A proposta original e as duas medidas que a redefiniram}
\label{cap:ancora}

A proposta que abriu esta linha de trabalho detectaria instalações
fotovoltaicas por imagem de satélite e cruzaria a detecção com o registro da
ANEEL, para verificar programas de crédito rural. Duas medidas a redefiniram
antes que qualquer modelo fosse treinado, e elas fecham caminhos distintos. A
primeira contamina o alvo de qualquer modelo de geração, inclusive de geração
centralizada. A segunda fecha o segmento que a proposta original queria
auditar.

\section{Magnitude e composição do corte no SIN}
\label{sec:fig1}

\begin{ficha}
  \item[Janela]      abril de 2024 a julho de 2026
  \item[Amostra]     85 usinas com apuração de \textit{constrained-off}
  \item[Unidade]     usina por intervalo de 30 minutos
  \item[Fonte]       ONS, restrição por \textit{constrained-off} fotovoltaico
                     e fator de capacidade
  \item[Teste]       nenhum: a análise é descritiva
\end{ficha}

O ONS publica, por intervalo restrito, a geração verificada e uma geração de
referência estimada a partir da curva de desempenho da própria usina. O
\termo{corte} é a diferença entre as duas, truncada em zero:

\[ C = \max\bigl(\text{referência} - \text{verificada},\; 0\bigr) \]

A definição vale apenas nos intervalos em que há razão de restrição apurada. O
campo \texttt{val\_\allowbreak geracaolimitada} não é a energia cortada. Ele registra o teto
ao qual a geração foi limitada. Sua correlação com a geração verificada é 0,86,
e com a perda apurada é 0,02. Usá-lo como energia cortada superestima a perda.

\begin{figure*}[t]
\centering
\includegraphics[width=\linewidth]{../figures/fig1_curtailment_sin.pdf}
\caption{\textbf{O corte de geração fotovoltaica no SIN é sistêmico, de razão
energética e incide sobre toda a janela solar.}
\textbf{a}, Perfil diurno médio da potência fotovoltaica das usinas com apuração
de \textit{constrained-off}. A área azul é a geração verificada e a área
vermelha é a energia cortada; a soma define a potência que estaria disponível.
Médias sobre todos os intervalos de 30 minutos do período, em horário de
Brasília, sem correção para hora solar local.
\textbf{b}, Taxa mensal de corte por subsistema, definida como energia cortada
dividida pela energia disponível.
\textbf{c}, Energia total cortada por razão de restrição, decomposta em origem
sistêmica e local. Os rótulos indicam o total por razão.
\textbf{d}, Distribuição espacial das usinas e conjuntos com apuração. A cor
codifica a taxa de corte acumulada e o tamanho, a capacidade instalada. A escala
de cor satura no percentil 97 para preservar contraste. Contornos estaduais do
IBGE.}
\label{fig:corte-sin}
\end{figure*}

A taxa de corte no parque com apuração é de 21,6\%. A razão energética responde
pela maior parte da energia cortada, e a origem declarada é predominantemente
sistêmica. O corte incide sobre toda a janela solar, e não sobre um horário de
ponta. As taxas do Nordeste e do Sudeste com Centro-Oeste crescem em ritmo
semelhante.

\conclusao{O corte é de razão energética e origem sistêmica, e não de falha
local. Isso o torna uma propriedade do despacho do sistema, e não um atributo
da usina cortada.}

\campo{Consequência para modelagem} Onde a geração de uma usina é limitada por
decisão operativa, um modelo que preveja geração a partir de irradiância ajusta
um alvo que incorpora essa decisão. O mesmo vale para um modelo que parta de
imagem de satélite. O erro não é aleatório, porque a decisão de cortar depende
de carga, de despacho térmico e de posição na rede.

\begin{objecao}
O denominador é o das usinas com apuração, e não o do parque solar nacional. A
geração distribuída tem regime de restrição distinto e não entra. A geração de
referência é estimativa do operador, e em 21,8\% dos intervalos restritos ela
fica abaixo da geração verificada; o truncamento em zero torna a estimativa
conservadora, e não exata. Razão e origem são classificações operativas do ONS,
e não inferência causal independente.
\end{objecao}

\campo{Comparação internacional} Os 21,6\% são de solar centralizada com
apuração. O ONS reporta 20,7\% para eólica e solar em 2025, com denominador
diferente e mesma ordem de grandeza. Os valores publicados para outros sistemas
são menores: 17\% da renovável variável disponível no Chile em 2024, cerca de
11\% na Espanha em julho de 2025, e 4,3\% de média para a solar da Califórnia.
O'Shaughnessy, Cruce e Xu \cite{oshaughnessy2020} organizam a série global de
corte fotovoltaico e mostram que a taxa cresce com a penetração e não com a
idade do sistema, o que põe o valor brasileiro na trajetória esperada e não fora
dela. Novan e Wang \cite{novan2024} medem que a taxa marginal é cerca de duas
vezes a média, o que implica exposição maior para uma usina que entre agora.

\section{Detectabilidade da geração distribuída}
\label{sec:fig2}

\begin{ficha}
  \item[Janela]      conexões registradas até abril de 2026
  \item[Amostra]     4,6 milhões de instalações
  \item[Unidade]     empreendimento
  \item[Fonte]       ANEEL, geração distribuída e informações técnicas
                     fotovoltaicas
  \item[Teste]       nenhum: a análise é descritiva sobre censo administrativo
\end{ficha}

A ANEEL publica a área do arranjo de cada instalação registrada. O número de
pixels que um sensor de resolução $r_s$ entrega sobre um arranjo de área
$A$ é

\[ \mathrm{px}(s) = A / r_s^2 \]

A conta transforma a resolução do sensor em uma grandeza comparável com o
registro, instalação por instalação.

\begin{figure*}[t]
\centering
\includegraphics[width=\linewidth]{../figures/fig2_mmgd_detectabilidade.pdf}
\caption{\textbf{A geração distribuída brasileira é majoritariamente sub-pixel
no Sentinel-2, e a série de conexões carrega uma descontinuidade regulatória.}
\textbf{a}, Distribuição acumulada da área do arranjo das instalações
registradas na ANEEL. As linhas tracejadas marcam a área de um pixel de cada
sensor, e o rótulo indica a fração de instalações menor que um pixel.
\textbf{b}, Fração de instalações que atinge o piso operacional de quatro
pixels, por classe de consumo e sensor.
\textbf{c}, Conexões mensais. A linha tracejada marca o fim da janela de
transição da Lei 14.300/2022, e o marcador indica o mês de pico.
\textbf{d}, Participação da classe rural no total de instalações de cada unidade
da federação. Contornos do IBGE.}
\label{fig:mmgd}
\end{figure*}

No Sentinel-2 a mediana é de 0,29 pixel por instalação, e a fração sub-pixel é
de 91,2\%. No CBERS-4A WPM são 7,25 pixels e 0,5\%. Num sensor comercial de
0,5 m são 116 pixels e fração próxima de zero. A classe rural, alvo dos
programas de crédito, tem área mediana de 41 m\textsuperscript{2} e 5,2\% de
segmentabilidade no Sentinel-2.

\conclusao{Nove em cada dez instalações de geração distribuída ocupam menos de
um pixel do Sentinel-2. A verificação por segmentação exige 2 m ou melhor.}

\campo{A resolução não é a restrição dominante} O registro não traz localização
utilizável. Os campos de coordenada existem no esquema e estão preenchidos em
0,06\% dos casos. Há o campo \texttt{CodCEP}, completo para os 7,7\% de pessoa
jurídica e truncado em cinco dígitos para os 92,3\% de pessoa física. A mediana
é de um prefixo por município, o que não refina o alvo rural. Sem coordenada por
instalação, a verificação exige varredura exaustiva do território. Apenas o
segmento de pessoa jurídica é endereçável.

\campo{A data de conexão não é exógena} A janela de transição da Lei
14.300/2022 concentrou 98 mil conexões em um único mês, quatro vezes a média do
ano anterior. Um estudo de evento sobre a data de conexão mediria a antecipação
regulatória, e não o efeito da política.

\begin{objecao}
Contar pixels não é o mesmo que detectar. Um objeto sub-pixel ainda altera a
assinatura espectral da célula. O que a contagem descarta é a segmentação, e não
a detecção espectral. O piso de quatro pixels é convenção declarada, e não
constante física. O conjunto é censo administrativo, e não amostra, de modo que
não há erro amostral a declarar.
\end{objecao}

\campo{Posição na literatura} Os trabalhos que segmentam telhados operam entre
0,1 e 0,3 m. O DeepSolar \cite{yu2018} usa imagem aérea submétrica, e o conjunto
multirresolução de Jiang et al. \cite{jiang2021} vai de 0,1 a 0,8 m. O
inventário global de Kruitwagen et al. \cite{kruitwagen2021} é a referência com
Sentinel-2, e cobre apenas instalações acima de 10 kW. Não foi localizado
trabalho que segmente telhado residencial a 10 m.
